% \autoref{sec:intro} identified two challenges in Connections. The first is lexical: recognising what four words have in common. The second is structural: all sixteen words must be partitioned into four disjoint groups at once, so a group cannot be judged on its own. This chapter is organised around that division.

% \autoref{sec:prior-work} reviews existing work on solving Connections, which has concentrated almost entirely on the first challenge. The remaining sections cover formalisms that apply to the second. \autoref{sec:puzzles-as-csp} describes how word puzzles are treated as weighted constraint satisfaction problems, which is how the partition requirement itself is expressed. \autoref{sec:information-theory} and \autoref{sec:strategies} turn to what happens when a puzzle is solved through repeated attempts rather than in one pass: how the information carried by feedback is measured, and how it is used to choose what to submit next.
\chapter{Related Work}
\label{chap:related_work}

This chapter provides an overview of prior work relevant to this thesis, including work that directly solves Connections using LLMs and work that inspired the methods used in this study. It first introduces prior work on solving Connections with AI. It then presents related work on the methods used, which serves as the knowledge base for this study.

\sectionnotoc{Prior Work on Connections Solving}
\label{sec:prior-work-connections}

Current work on this topic mainly uses LLMs as the main solution for solving Connections. These studies compare how well different models solve Connections. Some use Connections as a benchmark to test LLM capability, while others explore different prompting techniques and methods. \citeauthor{toddMissedConnectionsLateral2024} use three methods, namely a sentence-embedding baseline, a comparison between two LLMs (GPT-3.5-Turbo-1106 and GPT-4-1106-Preview), and \ac{CoT} prompting. Their experiments allow five incorrect and five invalid guesses per puzzle, which departs from the original game's rules. GPT-4 Turbo with \ac{CoT} achieves the highest success rate at 58.11\%. The best-performing sentence-embedding baseline (MPNet) reaches 11.6\%. LLMs are often stumped by categories that involve non-semantic word properties, abstract features, or context-dependent usage. They also report a finding relevant to this thesis. Puzzle order matters, since presenting puzzles in their original, grouped order achieves substantially higher accuracy than presenting them in randomized order~\cite{toddMissedConnectionsLateral2024}.

A similar work from \citeauthor{samadarshiConnectingDotsEvaluating2024} benchmarked several LLMs against both novice and expert human players and creates a taxonomy of knowledge types required to solve connections. They used in total 438 Connections puzzles and their best performing models Claude Sonnet 3.5 achieved 18\% success rate of fully solving the puzzles. They categorizes NYT Connections categories along two main dimensions, word form (covering morphology, orthography, phonology, and multiword expressions) and word meaning (covering semantic relations such as synonymy, polysemy, and hypernymy, along with associative and encyclopedic relations), with a further category for groups that combine both form-based and meaning-based cues. Results show that models performed comparably well on semantic relation categories but significantly worse on the other types.

More recent work from \citeauthor{loredolopezNYTConnectionsDeceptivelySimple2025} compare three game configurations: single attempt, four trys without feedback (one away prompt) and four trys with feedback. Human solvers outperform LLM under all three configurations. The best performing model Claude 3.5 under full hints reaches 40.25\%, while the human full-hints score reaches 60.67\%. They also experiment on different prompting techniques. \ac{CoT} and self-consistency prompting yield diminishing or even negative returns as puzzle difficulty increases. Simple input-output prompting sometimes outperforms these more elaborate reasoning strategies. Their herustic (non-LLM) baseline, using Multilingual-E5-Large-Instruct embedding and beam search, reaches 13.25\% accuracy.

\citeauthor{zhangLearningSemanticComplexities} compare in-context prompting of a large model (DBRX) against fine-tuning and distilling a smaller model (Mistral~7B) on solved puzzles. They found that fine-tuning a smaller model (Mistral 7B) directly on solved puzzles fails to capture the task itself, which underperformed even their clustering baseline. Their strongest result comes from Verbal Reinforcement Learning, in which a larger supervisor model critiques the fine-tuned model's output in natural language across three rounds, roughly doubling the average per-puzzle match rate relative to single-pass prompting \cite{zhangLearningSemanticComplexities}.

% \section{Lexical and Semantic Knowledge Resources}
% As discribed in chapter 1, this thesis also explores several forms of knowledge injection beyond raw embeddings. Those resources are mainly used to solve the semantic challenge of connections. They provide knowledge or explaination of the puzzle words and help form paulisble groupings of the puzzle. There are five sources examined in this thesis: WordNet, ConceptNet, Wikipedia, Datamuse, Language model (Mistral small 24b). They differ not only in coverage but in the nature of the knowledge they provide. WordNet and ConceptNet organize knowledge as explicit relations between words or concepts, while Wikipedia provides unstructured encyclopedic text that must be mined for relevant context. Datamuse offers a narrower, more mechanical form of lexical association, surfacing rhymes, phonetic near-matches, and statistically co-occurring words rather than just semantic relations. And lastly, we prompt a large language model to generate meaning for each word on demand, an approach related to the NLP task of definition modeling (Noraset et al., 2017), in which a model produces a dictionary-style definition given a word or its embedding. Unlike the four resources above, the knowledge here is not stored in an inspectable structure but generated parametrically at query time, and may vary across generations of the same word.

% Two knowledge resources recur across the enhancement methods discussed above: WordNet and ConceptNet. WordNet (Miller, 1995) organizes English vocabulary into synsets, groups of words sharing a specific sense, with distinct synsets maintained for each sense of a polysemous word. This structure is what makes sense-level, rather than word-level, embedding possible in principle. ConceptNet (Speer et al., 2017) instead represents common-sense knowledge as a graph of concepts connected by typed relations (e.g., "is used for," "is a kind of"), offering broader coverage of everyday associative knowledge than WordNet's more strictly lexical relations, but with correspondingly weaker coverage of proper nouns and domain-specific or pop-culture references, an asymmetry noted in Section [X].

% The limitation that motivates moving beyond simple word-level embeddings in the first place is well documented in the NLP literature under the term meaning conflation deficiency: representing every sense of a word with a single vector causes distinct, sometimes unrelated meanings to be averaged into one representation, degrading the precision of that representation for any individual sense (Schütze, 1998; Camacho-Collados and Pilehvar, 2018). This deficiency is not specific to any one embedding model or corpus; it is a structural property of representing word types, rather than word senses, as points in a vector space. As we show in Section [X], this deficiency accounts for the majority of scoring failures observed in this thesis, motivating a sense-level treatment of word similarity rather than a search for additional knowledge sources per se.
% Wikipedia has been widely used in NLP as a general-purpose source of encyclopedic and named-entity knowledge, owing to its broad, crowd-maintained coverage of proper nouns, cultural references, and domain-specific topics that lexical resources such as WordNet do not systematically include. Datamuse, by contrast, is a public API for retrieving words related to a query along narrow, mechanically defined axes such as rhyme, spelling similarity, or corpus co-occurrence, rather than curated semantic content.

% Generating word glosses with a language model, rather than retrieving them from a fixed resource, raises a different question from the ones addressed by WordNet, ConceptNet, or Wikipedia: not whether the relevant knowledge exists in some accessible structure, but whether a generative model can supply it directly, without requiring that structure to exist at all. This distinction matters for the argument developed in this thesis. As shown in Section [6.3], however, the knowledge that proves systematically missing (figurative senses and proper-noun or pop-culture associations) is absent regardless of whether the source is a fixed lexical resource, a knowledge graph, unstructured encyclopedic text, or a language model generating glosses on demand. The failure recurs across sources that differ fundamentally in how their knowledge is stored and accessed, suggesting that the bottleneck lies not in any particular resource's coverage, but in how that knowledge, however obtained, is applied: at the level of isolated candidate groups rather than the global partition (Section [X]).

% wordnet
% conceptnet
% wikipedia
% datamuse
% embedding
% llm
% 两者的区别

% Related Work（真正意义上的）：回顾"别人做过什么、和你的工作是什么关系"——你需要引用别人的研究结论/方法，用来对比、定位自己的贡献。比如 Todd/Samdarshi/Lopez（别人怎么测这个任务）、Knuth/Wordle（别人怎么解决同类交互式问题）、CSP先例（别人怎么形式化同类谜题）——这些都是"别人做了什么、我在哪里接续/不同"，是比较性的内容，理应留在Related Work。

% Background/Preliminaries（背景知识）：读者看懂你的Method需要的术语和工具介绍——WordNet的synset结构是什么、ConceptNet的关系图是什么、义项稀释是什么现象。这些不是"别人的研究和你的关系"，而是"你用到的工具本身是什么"。这类内容更适合放在你实际使用它的地方，而不是提前在Related Work里打包介绍一遍让读者硬记。

% 按这个标准，"知识资源"这一节里的内容其实应该拆开：

% WordNet的synset结构、义项稀释(meaning conflation deficiency)——这是你4.2节MaxSynset方法的直接理论依据，应该就地放进第4章方法论，读者读到"为什么这样设计MaxSynset"时立刻看到支撑引用，不需要翻回前面Related Work去找。
% ConceptNet、Wikipedia、Datamuse的介绍——这三个是第6章"知识增强否定链"里各自实验用到的资源，应该放在第6章对应实验小节的开头，各自用一两句话介绍完就直接讲实验设计和结果，不需要在Related Work里单独占一节。
% 调整后 Related Work 的实际内容（回到只保留"真正比较性"的三节）
% Prior Work on Connections Solving（Todd/Samdarshi/Merino/Lopez/Zhang）
% 谜题作为约束满足问题（CSP/Dr.Fill先例）
% Sequential Decision-Making Under Partial Feedback（Knuth/Wordle/group testing）

% "知识资源"这一节直接取消，内容拆分到第4章和第6章对应位置去，各自轻量地引用（比如第4章开头一两句带上Schütze/Camacho-Collados的引用，第6章每个实验小节开头一两句带上WordNet/ConceptNet/Wikipedia各自的出处）。

% 这样调整的好处：Related Work只剩"和别人比较"的内容，读起来更聚焦；背景知识分散到各自被用到的地方，读者不需要提前记住一堆定义再等好几页才看到它们被用上。

% 你觉得这个拆分方式可以吗？如果同意，我可以先帮你把 Method 第4.2节 补上 Schütze/Camacho-Collados 的引用（之前写的时候还没有这条文献），第6章各实验小节开头的资源介绍等你写到那里时再处理。

\sectionnotoc{Word Puzzles as Weighted \ac{CSP}}
\label{sec:puzzles-as-csp}

% what is csp
\ac{CSP}s are mathematical questions defined as a set of objects whose state must satisfy a number of constraints or limitations \cite{ConstraintSatisfactionProblem2026}. A (finite domain) \ac{CSP} can be expressed in the following form: given a set of variables, together with a finite set of possible values that can be assigned to each variable, and a list of constraints, find values of the variables that satisfy every constraint \cite{brailsfordConstraintSatisfactionProblems1999}. It is formally defined as a triple $\langle X, D, C \rangle$ where X is a set of variables, D specifies a domain of possible values for each variable, and C is a set of constraints restricting which combinations of value assignments are permitted \cite{ConstraintSatisfactionProblem2026}.

% what is weighted csp
In artificial intelligence, a Weighted Constraint Satisfaction Problem, also known as Valued Constraint Satisfaction Problem, is a constraint satisfaction problem in which preferences among solutions can be expressed \cite{gallardoSolvingWeightedConstraint2009}. A plain \ac{CSP}, treats every solution as equally good: it either satisfies the constraints or it does not. Connections needs more than that. Not every four-groups-of-four split is a correct solution, only one is. What is available instead is a measure of how well four words seem to belong together, based on their similarity. Thus, Connections can be defined as a weighted \ac{CSP}. It keeps the hard constraint (a valid four-groups-of-four partition) but adds a score to each candidate value assignment, so that solving the problem means finding the constraint-satisfying assignment with the best total score, not just any assignment that satisfies the constraints.

% dr fill and related literature
\subsectionnotoc{Solving Word Puzzles as Weighted \ac{CSP}}
Applying this concept to word puzzles is not new. \citeauthor{ginsbergDrFillCrosswordsImplemented2014} formalizes American-style crossword solving as a weighted \ac{CSP} in the program Dr\.Fill: each grid cell is a variable, each candidate fill is a value, and the constraints require intersecting words to agree on shared letters \cite{ginsbergDrFillCrosswordsImplemented2014}. Because many fills are plausible for a given clue, each candidate fill is given a likelihood-based weight, and Dr.Fill searches for the assignment that maximizes total weight while still satisfying the grid's hard constraints. It draws candidate fills from several sources: a database of over 47,000 previously published crosswords with 1.9 million unique clues, two dictionaries of varying size, 1.2 million synonyms drawn from WordNet and other online sources, and a database of Wikipedia titles and word pairs. Each candidate fill is then scored on five criteria: whether the clue matches one seen before in the clue database, whether the fill's part of speech agrees with the part of speech inferred from the clue, the fill's crossword "merit" (a hand-rated measure of how desirable a word is as a crossword answer in general), whether the fill is an abbreviation, and whether the clue is a fill-in-the-blank type resolvable against Wikipedia text. 
% These five criteria are combined linearly, with the weights tuned on a fixed testbed of 100 puzzles to maximize the number of words entered correctly before the first mistake. 
Dr.Fill then searches for the assignment that maximizes total score while still satisfying the grid's hard letter-intersection constraints. Two further design choices in Dr.Fill are worth noting. First, they found that looking further ahead before committing to a choice gave no real benefit, and the system used only a single step of lookahead instead. Second, branch-and-bound, a standard pruning technique for weighted CSPs, was tested and dropped because it added little value in this domain, as a result postprocessing was used instead. In both cases, a more sophisticated search method did not beat a simpler one, once the scoring function and hard constraints were doing most of the work. Since complete knowledge of what each clue means is impractical to obtain, Dr.Fill relies mainly on the hard constraint that crossing words must share letters to narrow down candidates, not on understanding the clues themselves. Dr.Fill inspired the methodology used in this thesis, which also treats Connections as a weighted \ac{CSP} and uses a scoring function to evaluate candidate groupings.

\subsectionnotoc{Set Partitioning}
The constraint $C$ of Connections (every word belongs to exactly one group and every group contains exactly four words), more precisely speaking, is a combinatorial optimization problem known as set partitioning. A partition of a set is a way of dividing all of its elements into non-empty, pairwise disjoint subsets whose union is the original set \cite{PartitionSet2026}. The set partitioning problem asks, given a ground set and a collection of candidate subsets each associated with a cost, to select a subset of these candidates that forms a partition of the ground set at minimum or maximum total cost \cite{muellerSolvingSetPartitioning1998}. The sixteen Connections puzzle words are the ground set, candidate four-word groups are subsets among which a selection must be made, and each candidate group's score from the weighted \ac{CSP} explained above serving is its cost. Identifying this structure is what motivates the grouping strategy adopted later in this thesis.

Formulating a puzzle as a weighted CSP defines what to search for. Among all groupings that satisfy the hard constraint, the goal is the one with the highest total score. This does not guarantee that the grouping is correct. The hard constraint only rules out invalid partitions. Many valid partitions remain, and the scoring function is what distinguishes them from each other. If the scoring function is imperfect, the highest-scoring grouping can still be wrong. Puzzle games like Connections give players a further source of information, namely the response the game returns after each attempt. Using this response turns the problem from a single optimization into a sequence of attempts, where each attempt narrows down what remains possible. \autoref{sec:information-theory} covers how much a single response is worth, and \autoref{sec:strategies} covers what to do with it.

% wordle: infromation theory
\sectionnotoc{Information Theory}
\label{sec:information-theory}
% % what is information theory, entropy, and information gain
% Information theory is the mathematical study of the quantification, storage, and communication of a particular type of mathematically defined information \cite{InformationTheory2026}. One of its central concepts is entropy, which measures the uncertainty associated with a random variable. In the context of puzzle solving, entropy can be used to quantify the amount of information gained from a particular guess or action. The concept of information gain is closely related, representing the reduction in uncertainty achieved by making a specific choice or observation.
% % wordle

% information theory and entropy
Information theory is the mathematical study of the quantification, storage, and communication of a particular type of mathematically defined information \cite{InformationTheory2026}. It gives information a precise mathematical definition grounded in probability \cite{shannonMathematicalTheoryCommunication1948}. Gaining information and reducing uncertainty about a random event are the same thing. By convention, information is measured in base-2 logarithms, giving the bit as its unit: each time half of the probability mass is ruled out, exactly one bit is gained \cite{pinkardVisualIntroductionInformation2026}. The information contained in an outcome with probability $p(x)$ is $\log_2 \frac{1}{p(x)}$.

The probability-weighted average of this information across all outcomes is the entropy, denoted $H(\mathrm{X})$:
\begin{equation}
H(\mathrm{X}) = \sum_{x \in \mathcal{X}} p(x) \log_2 \frac{1}{p(x)}.
\label{eq:entropy}
\end{equation}
Entropy serves as the connection between probability and information content. It measures how uncertain a random event is on average. Higher entropy means more uncertainty to begin with, so more information must be acquired to become certain about the outcome. Entropy is maximized when probability is spread evenly over all outcomes. In that case it reduces to
\begin{equation}
H_{\text{max}}(\mathcal{X}) = \log_2 |\mathcal{X}|,
\label{eq:maxentropy}
\end{equation}
where $|\mathcal{X}|$ is the number of possible outcomes \cite{pinkardVisualIntroductionInformation2026}.

Conditional entropy, $H(\mathrm{X} \mid y)$, measures how much uncertainty about $\mathrm{X}$ remains after observing an outcome $y$. The reduction in entropy achieved by an observation is the information gain. Taken in expectation over all possible outcomes of an observation, before that observation is actually made, this gives the expected information gain of an action:
\begin{equation}
IG(\text{guess}) = H(\mathrm{X}) - \mathbb{E}\big[H(\mathrm{X} \mid \text{feedback})\big].
\label{eq:infogain}
\end{equation}
This quantity is what solvers use to decide which action is most useful before knowing what the feedback will be. An action with higher expected information gain is expected to narrow down the remaining candidates more, regardless of which specific feedback is eventually returned.

\subsectionnotoc{Solving Wordle with Information Theory}
% wordle
Wordle is another word puzzle game from NY Times. It asks the player to guess a five-letter word within six attempts. After each guess, every letter is marked as correct and in the right position, present but in the wrong position, or absent, giving detailed per-letter feedback. Information theory based Wordle solvers apply this framework directly: at each step, the candidate word maximizing expected information gain is chosen as the next guess \cite{aladailehSolvingWordleUsing2026}. According to Grant Sanderson pointed out in his video, the best possible first guess yields an expected information gain of 5.89 bits, reducing the remaining uncertainty by roughly half in a single observation \cite{3blue1brownSolvingWordleUsing2022}. \citeauthor{aladailehSolvingWordleUsing2026} report that entropy-based selection outperforms a heuristic that selects words by letter distribution, and propose the approach as a general baseline for applying Shannon entropy to other games \cite{aladailehSolvingWordleUsing2026}.
% connections
Information theory can be applied to Connections as well. Detailed description see \autoref{chap:problem-formulation}.
% Before any guess is made, the number of possible partitions of the sixteen words defines the entropy of the hypothesis space, in the same way that the set of allowed words defines the entropy of the candidate space in Wordle. Each guess, a submitted group of four words, returns feedback, and this feedback can in principle be used the same way: to compute an expected information gain and to eliminate partitions inconsistent with the observed outcome, reducing entropy at each step.

% The key difference is the feedback signal itself. A Wordle guess returns a detailed, per-letter pattern, one of many possible outcomes, and consequently a large expected information gain. A Connections guess returns only one of three coarse outcomes: correct, one away, or wrong. By \autoref{eq:maxentropy}, this limits the maximum information a single guess can carry to $\log_2 3 \approx 1.58$ bits, regardless of which four words are guessed or how the guess is chosen. Over the four guesses allowed by the game's real rules, this places a hard upper bound of roughly 6.3 bits on the total information obtainable through feedback alone, far below the entropy of the hypothesis space itself. This bound is developed with exact numbers in Section~\ref{sec:problem-formulation}, and is the central reason, developed later in this thesis, why decision-strategy sophistication cannot substitute for the quality of the prior belief formed before any guess is made.

\sectionnotoc{Strategies for Puzzle Games}
% 前面几节已经建立了两个身份：CSP/set partitioning（结构性的、静态的）、information theory（工具性的）。
\label{sec:strategies}
\autoref{sec:puzzles-as-csp} introduces the structural constraint, and then \autoref{sec:information-theory} as a way of measuring how much a single observation can reveal. After measuring information, a solver needs a rule that turns the current state of knowledge into the next move. This section covers the other question: given feedback from previous attempts, how should the next guess be chosen?

% \citeauthor{cunananOptimalStrategiesWordle2023} address this question by defining a general class of games to which Wordle belongs \cite{cunananOptimalStrategiesWordle2023}. A guessing game is a tuple $(G, S, R, r^*, a)$, where $G$ is the set of allowable guesses, $S$ the set of possible secrets, $R$ the set of possible responses, $r^*$ the response indicating the secret has been found, and $a: G \times S \to R$ the answering function that determines which response a guess receives. After a sequence of guesses and responses, the candidate set is the set of secrets consistent with every response observed so far, and a strategy is a rule mapping each candidate set to the next guess to submit.

% Within this framework, strategies are expressed as valuations: functions that score each possible guess against the current candidate set. \citeauthor{cunananOptimalStrategiesWordle2023} collect several valuations previously developed for Mastermind, including minimising the size of the largest remaining candidate set, maximising information gain in the sense of the previous section, and maximising the number of distinct responses a guess can produce. Evaluating these on Wordle, they find that none performs best alone, and that maximising information gain in particular does not yield the optimal strategy \cite{cunananOptimalStrategiesWordle2023}.

\citeauthor{cunananOptimalStrategiesWordle2023} provide a solution for Wordle and for a wider class of games with the same structure \cite{cunananOptimalStrategiesWordle2023}. After a sequence of guesses and responses, the candidate set is the set of possible secrets that remain consistent with every response observed so far. A strategy is a rule that maps the current candidate set to the next guess, and different strategies can be compared by scoring every possible guess against the candidate set.

Several such scoring functions were developed earlier for Mastermind. These functions include minimizing the size of the largest candidate set that could remain, maximizing information gain as described in the previous section, and maximizing the number of distinct responses a guess can produce ~\cite{knuthComputerMasterMind1976}. \citeauthor{cunananOptimalStrategiesWordle2023} test these functions on Wordle and find that none performs best on its own. Maximizing information gain in particular does not yield the optimal strategy ~\cite{cunananOptimalStrategiesWordle2023}.